\section{Dataset Description}\label{sec:dataset-description}
The present work utilized ultra-long-term \gls{sqeeg} data acquired during two clinical trials using the \textit{24/7 EEG™ SubQ system} (UNEEG Medical A/S, Denmark). Consequently, two distinct datasets were analyzed. 

\subsection{Recording device}
The system consists of an implantable component with three electrodes positioned over the presumed seizure onset zone and an external recording unit, resulting in two-channel recordings.
The implantation hemisphere and electrode placement were determined individually based on each patient’s epileptic focus.
Data were recorded exclusively when the external unit was connected to the implant.

\begin{figure}[h]
    \centering

    \begin{subfigure}[t]{0.48\textwidth}
        \centering
        \includegraphics[width=\textwidth]{UNEEG_implant}
        \caption{Implant and electrodes.}
        \label{fig:uneeg-implant}
    \end{subfigure}
    \hfill
    \begin{subfigure}[t]{0.48\textwidth}
        \centering
        \includegraphics[width=\textwidth]{UNEEG_device_orientation}
        \caption{External recording unit orientation.}
        \label{fig:uneeg-orientation}
    \end{subfigure}

    \caption[24/7 EEG\texttrademark{} SubQ system components.]{24/7 EEG\texttrademark{} SubQ system components. Source: \textcite{uneeg_medical_uneeg_nodate}}
    \label{fig:uneeg-system}
\end{figure}


\subsection{Dataset}
The first dataset was obtained at the Epilepsy Center Freiburg, Germany, as part of the multicenter ULTRA2 study.
After implantation, patients were instructed on how to use the device, and concurrent video-\gls{eeg} and subcutaneous \gls{eeg} recordings were obtained during the inpatient period in the \gls{emu}. 
Upon discharge, patients continued with outpatient recordings using the device. 
Approximately monthly visits to the clinic were conducted for follow-up and data transfer throughout the individual study duration. 

The seizure annotations were created using a semi-automatic scheme. 
First, an automated classifier \parencite{furbass_automatic_2017} processed the \gls{eeg} recordings to detect potential seizures. 
These detections were then reviewed by two epileptologists to verify true seizure events. 
In cases of disagreement, a third expert made the final decision. 
The annotations consisted of a single marker per event, which indicated the seizure onset detected by the classifier, whereas detailed information on seizure onset, offset, and type was only sparsely available. 

The second dataset comprised recordings from three participants enrolled in the SUBER study (“SUBcutaneous EEG: foRecasting of epileptic seizures through
investigation of long-term dynamics of seizure occurrences, stress, sleep, and other factors”)---an observational, prospective, non-randomized, study \parencite{viana_real-world_2025}, conducted at King’s College London and King’s College Hospital NHS Foundation Trust.  
Their data were shared as part of a worldwide seizure prediction competition. 

Both datasets initially consisted of seizure annotations and \gls{eeg} recordings on a per-patient basis. 
Patients are subsequently referred to by the initial letter of their respective dataset (\underline{U}LTRA2/\underline{C}ompetition), followed by their ID number, e.g., U17 or C03.

Both studies were conducted in accordance with the Declaration of Helsinki and approved by the respective local ethics committees (20-1175; 19/LO/0354). 
Informed consent was obtained from all participants.
